% !Mode:: "TeX:UTF-8"

\chapter{基于Feign的声明式服务调用}

\section{设计原理}

Feign是Netflix开发的声明式HTTP客户端，通过"接口 + 注解"的编程模型大幅简化了服务间调用代码。其核心优势包括：

\begin{itemize}
	\item \textbf{声明式编程模型}：只需定义Java接口并标注注解，Feign自动生成代理实现类
	\item \textbf{内置负载均衡}：集成Spring Cloud LoadBalancer，自动实现客户端负载均衡
	\item \textbf{集成熔断降级}：通过\verb|fallback|属性指定降级实现类
	\item \textbf{契约驱动}：Feign接口定义本身就是服务调用契约
\end{itemize}

\section{Feign客户端定义}

项目中共定义了9个Feign客户端，涵盖所有跨服务调用场景：

\begin{table}[H]
	\centering
	\caption{Feign客户端总览}
	\begin{tabular}{p{3.5cm}p{3.0cm}p{4.0cm}}
		\hline
		\textbf{客户端} & \textbf{所在服务} & \textbf{目标服务与用途} \\
		\hline
		BusinessClient & food-server & business-server：验证商家存在 \\
		FoodFeignClient & cart-server & food-server：获取食品信息 \\
		OrderFeignClient & cart-server & business-server：创建订单 \\
		BusinessClient & order-server & business-server：获取商家信息 \\
		CartClient & order-server & cart-server：购物车操作 \\
		PointClient & order-server & point-server：积分操作 \\
		VirtualWalletClient & order-server & wallet-server：余额、转账 \\
		UserClient & wallet-server & user-server：验证用户存在 \\
		OrderClient & point-server & order-server：关联订单积分 \\
		\hline
	\end{tabular}
\end{table}

以下以订单服务调用钱包服务为例：

\begin{lstlisting}[language=Java, caption=VirtualWalletClient.java]
@FeignClient(name = "elm-wallet-server")
public interface VirtualWalletClient {
    @GetMapping("/wallet/getBalance")
    CommonResult<BigDecimal> getBalanceByUserId(
        @RequestParam("userId") String userId);

    @PostMapping("/wallet/transfer")
    CommonResult<Integer> transfer(
        @RequestParam("fromUserId") String fromUserId,
        @RequestParam("toUserId") String toUserId,
        @RequestParam("amount") BigDecimal amount);
}
\end{lstlisting}

\section{Feign集成Fallback降级}

当目标服务不可用时，Feign通过\verb|fallback|属性调用降级实现类，返回兜底数据：

\begin{lstlisting}[language=Java, caption=Feign Fallback降级]
@FeignClient(name = "elm-business-server",
             fallback = BusinessClientFallback.class)
public interface BusinessClient {
    @GetMapping("/businesses/{businessId}")
    String getBusinessById(@PathVariable("businessId") Integer businessId);
}

@Component
public class BusinessClientFallback implements BusinessClient {
    @Override
    public String getBusinessById(Integer businessId) {
        return "{\"code\":200, " +
               "\"message\":\"Fallback: Business service unavailable\"}";
    }
}
\end{lstlisting}

食品服务调用商家服务验证商家是否存在时，若商家服务不可用，Fallback返回模拟的成功响应，使食品服务仍能正常运行（宽松策略，假设商家存在）。这种设计体现了微服务架构中"优雅降级"的思想。

\section{组合调用示例}

订单服务创建订单时，组合使用多个Feign客户端：

\begin{lstlisting}[language=Java, caption=订单创建中的Feign组合调用]
@Transactional
public int createOrders(Orders orders) {
    // 1. Feign获取购物车
    CommonResult<?> cartResult =
        cartClient.listCart(orders.getUserId(), orders.getBusinessId());
    // 2. 创建订单主记录与明细
    ordersMapper.saveOrders(orders);
    orderDetailetMapper.saveOrderDetailetBatch(list);
    // 3. Feign清空购物车
    cartClient.removeCart(orders.getUserId(), orders.getBusinessId());
    return orderId;
}
\end{lstlisting}

整个过程是声明式的——开发者只需关注业务逻辑，HTTP通信的细节完全由Feign处理。
